\subsection{Diagonalization of Hilbert-Schmidt mappings}

THEOREM 2. --- Let E and F be two hilbertian spaces and u a Hilbert-Schmidt mapping from E into F. There exists an orthonormal basis \((e_{i})_{i\in I}\) of E which is transformed by u into an orthogonal family in F.

Let B denote the (closed) unit ball of E, with the weakened topology assigned to it; this is a compact space (V, p.~17). We put \(\mathbf{Q}(x) = \|u(x)\|^{2}\) for all \(x \in B\) . Finally let P denote the set of all vectors x in E satisfying the following property;

\begin{enumerate}
\def\labelenumi{(\Alph{enumi})}
\setcounter{enumi}{7}
\tightlist
\item
  For every \(y \in \mathbf{E}\) orthogonal to \(x\) , the element \(u(y)\) of \(\mathbf{E}\) is orthogonal to \(u(x)\) .
\end{enumerate}

Lemma 4. --- The function \(\mathbf{Q}:\mathbf{B}\to \mathbf{R}\) is continuous.

Let \((f_j)_{j\in \mathbf{J}}\) be an orthonormal basis of F. Put \(\lambda_{j} = \| u^{*}(f_{j})\|^{2}\) for all \(j\in \mathbf{J}\) . Since \(u\) belongs to \(\mathcal{L}^2 (\mathrm{E};\mathrm{F})\) we have \(u^{*}\in \mathcal{L}^{2}(\mathrm{F};\mathrm{E})\) , hence \(\sum_{j}\lambda_{j} < + \infty\) . Further, we have

\[
Q (x) = \left\| u (x) \right\| ^ {2} = \sum_ {j} \left| \langle u ^ {*} (f _ {j}) | x \rangle \right| ^ {2}\tag{43}
\]

by Parseval's formula (V, p.~22) and the definition of the adjoint (V, p.~38). For every \(x \in \mathbf{B}\) , \(|\langle u^{*}(f_{j})|x\rangle|^{2} \leqslant \lambda_{j}\) by Cauchy-Schwarz inequality; consequently, the convergence of the sum in formula (43) is uniform on B, hence lemma 4 (GT, X, § 1, No.~6).

Lemma 5. --- Let \(\mathbf{E}_1\) be a closed vector subspace of \(\mathbf{E}\) , stable under \(u^*u\) . If \(\mathbf{E}_1 \neq \{0\}\) , then there exists a vector of norm 1 in \(\mathbf{E}_1 \cap \mathbf{P}\) .

Since B is weakly compact, so is the weakly closed subspace \(B \cap E_{1}\) of B. Hence there exists (GT, IV, §6, No.~1, th. 1) a point \(x_{0}\) in \(B \cap E_{1}\) such that \(Q(x_{0} \geqslant Q(x))\) for all \(x \in B \cap E_{1}\) . If \(Q(x_{0}) = 0\) , we have \(Q(x) = 0\) and so \(u(x) = 0\) for all \(x \in B \cap E_{1}\) . Thus \(E_{1} \subset P\) and lemma 5 follows in this case.

Suppose now that \(\mathbf{Q}(x_0) > 0\) , then \(x_0 \neq 0\) . Since the vector \(\| x_0 \|^{-1} \cdot x_0\) belongs to \(\mathbf{B} \cap \mathbf{E}_1\) , we have

\[
\mathrm{Q} \left(x _ {0}\right) \geqslant \mathrm{Q} \left(\left\| x _ {0} \right\| ^ {- 1}. x _ {0}\right) = \mathrm{Q} \left(x _ {0}\right) / \left\| x _ {0} \right\| ^ {2}
\]

i.e.~\(\|x_{0}\|=1\) . We shall prove that \(x_{0}\) belongs to P; let \(y\in E\) be orthogonal to \(x_{0}\) . It is enough to prove that \(u(y)\) is orthogonal to \(u(x_{0})\) . But since y is the sum of a vector of \(E_{1}\) and a vector orthogonal to \(E_{1}\) , and both orthogonal to \(x_{0}\) (since \(x_{0}\in E_{1}\) ), it is enough to consider the following two cases:

\begin{enumerate}
\def\labelenumi{\alph{enumi})}
\item
  \(y\) is orthogonal to \(\mathbf{E}_1\) : since \(\mathbf{E}_1\) is stable under \(u^* u, u^* u(x_0) \in \mathbf{E}_1\) , hence \(0 = \langle y | u^* u(x_0) \rangle = \langle u(y) | u(x_0) \rangle\) .
\item
  \(y\) belongs to \(\mathbf{E}_1\) : for all \(t \in \mathbf{R}\) , the vector \(x(t) = (x_0 + ty) / \| x_0 + ty \|\) belongs to \(\mathbf{B} \cap \mathbf{E}_1\) . We have \(Q(xt) = f(t) / g(t)\) with
\end{enumerate}

\[
f (t) = \left\| u (x _ {0}) \right\| ^ {2} + 2 t \mathscr {R} \left\langle u (x _ {0}) | u (y) \right\rangle + t ^ {2} \left\| u (y) \right\| ^ {2}
\]

\[
g (t) = 1 + t ^ {2} \| y \| ^ {2}.
\]

In view of the definition of \(x_0\) , we have \(\mathrm{Q}(x(0)) \geqslant \mathrm{Q}(x(t))\) for all real \(t\) , hence \(\frac{d}{dt} \mathrm{Q}(x(t))\) is zero for \(t = 0\) . But \(f(0) = \| u(x_0)\|^2\) , \(g(0) = 1\) , \(f'(0) = 2\mathcal{R}\langle u(x_0)|u(y)\rangle\) , \(g'(0) = 0\) . Since

\[
\frac {d}{d t} \mathrm{Q} (x (t)) = \frac {f ^ {\prime} (t) g (t) - f (t) g ^ {\prime} (t)}{g (t) ^ {2}},
\]

we conclude that \(f'(0)=0\) , that is, \(\mathcal{R}\langle u(x_{0})|u(y)\rangle=0\) . When K=R, \(u(x_{0})\) is orthogonal to \(u(y)\) , when K=C, the vector iy belongs to \(E_{1}\) and is orthogonal to \(x_{0}\) , hence \(\mathcal{I}\langle u(x_{0})|u(y)\rangle=-\mathcal{R}\langle u(x_{0})|u(iy)\rangle=0\) , and finally \(u(x_{0})\) is orthogonal to \(u(y)\) . This proves lemma 5.

Now we prove th. 2. Applying th. 1 of S, III, §4, No.~5 we see, as in V, p.~23, that there exists a set S which is maximal among the orthonormal subsets of E contained in P. Let \(E_{1}\) be the set of all vectors orthogonal to S. Let \(y \in E_{1}\) ; if \(x \in S\) , the vectors x and y are orthogonal, and since \(S \subset P\) , we conclude that \(u(x)\) and \(u(y)\) are orthogonal; then

\[
\langle x | u ^ {*} u (y) \rangle = \langle u (x) | u (y) \rangle = 0
\]

and \(u^{*}u(y)\) is orthogonal to S. Hence \(E_{1}\) is stable under \(u^{*}u\) . If we had \(E_{1} \neq \{0\}\) , there would exist a vector \(x\) of norm 1 in \(E_{1} \cap P\) (lemma 5) and \(S \cup \{x\}\) would be an orthonormal subset of E contained in P. This would contradict the maximal character of S. Hence \(E_{1} = \{0\}\) and S is an orthonormal basis of E. Q.E.D.

COROLLARY 1. --- Let v be a continuous, positive endomorphism with finite trace of the hilbertian space E. There exists an orthonormal basis \((e_{i})_{i\in I}\) of E and a summable family of positive real numbers \((\lambda_{i})_{i\in I}\) such that \(v(e_{i}) = \lambda_{i}e_{i}\) for all \(i \in I\) .

Let \(\Phi(x, y) = \langle x | v(y) \rangle\) for \(x, y\) in E. Then \(\Phi\) is a positive hermitian form on E. There exists (V, p.~8, corollary) a hilbertian space F and a continuous linear mapping \(u\) from E into F such that \(\Phi(x, y) = \langle u(x) | u(y) \rangle\) for \(x, y\) in E. In other words, we have \(v = u^{*}u\) . By virtue of def. 9 (V, p.~52), \(u\) is a Hilbert-Schmidt mapping from E into F. By th. 2, there exists an orthonormal basis \((e_{i})_{i \in I}\) of E such that the vectors \(u(e_{i})\) are two by two orthogonal. Let \(i \in I\) ; for every \(j \neq i\) in I, we have

\[
\langle e _ {j} | v (e _ {i}) \rangle = \langle u (e _ {j}) | u (e _ {i}) \rangle = 0
\]

hence \(v(e_{i})\) is proportional to \(e_{i}\) and is of the form \(\lambda_{i}e_{i}\) , where \(\lambda_{i}=\langle e_{i}|v(e_{i})\rangle\) ; then

\[
\lambda_ {i} \geqslant 0 \quad \text { and } \quad \sum_ {i \in I} \lambda_ {i} = \operatorname{Tr} (v) <   + \infty .
\]

COROLLARY 2. --- Let E be a hilbertian space. Then \(\mathcal{L}^{1}(\mathrm{E})\subset \mathcal{L}^{2}(\mathrm{E})\)

The real case reduces to the complex case by the extension of scalars; we can therefore assume that \(\mathbf{K} = \mathbf{C}\) .

Since \(\mathcal{L}^{2}(E)\) is a vector subspace of \(\mathcal{L}(E)\) , it is enough to prove that every continuous and positive endomorphism v of E with finite trace belongs to \(\mathcal{L}^{2}(E)\) . With the notations of cor. 1, we have

\[
\sum_ {i \in \mathbf {I}} \left\| v (e _ {i}) \right\| ^ {2} = \sum_ {i \in \mathbf {I}} \lambda_ {i} ^ {2} \leqslant (\sum_ {i} \lambda_ {i}) ^ {2} <   + \infty .
\]

COROLLARY 3. --- Let v be a continuous positive endomorphism of the hilbertian space E with a finite trace. There exists a positive Hilbert-Schmidt endomorphism w of E such that \(v = w^{2}\) and such that v commutes with w.

With the notations of cor. 1, it is enough to consider the endomorphism \(w\) which transforms the vector \(\sum_{i\in I}\xi_i e_i\) into the vector \(\sum_{i}\lambda_i^{1 / 2}\xi_i e_i\) .

Remark. --- With the notations of th. 2, let J be the set of all \(i \in \mathbf{I}\) such that \(u(e_i) \neq 0\) . For all \(i \in \mathbf{J}\) , let \(\lambda_i = \| u(e_i)\|\) and \(f_i = \lambda_i^{-1}u(e_i)\) . Then \((e_i)_{i \in \mathbf{J}}(\text{resp. } (f_i)_{i \in \mathbf{J})}\) is an orthonormal basis of the initial (resp. final) subspace of \(u\) , we have \(u(e_i) = \lambda_i f_i\) for all \(i \in \mathbf{J}\) and \(\sum_{i \in \mathbf{J}} \lambda_i^2 = \| u\|_2^2\) is finite.

